\documentclass[11pt, a4paper]{article}

% bfw-style.tex — gemeinsame Style-Definitionen für alle Handouts/Aufgabenhefte
% Voraussetzung: Kompilation mit xelatex (wegen fontspec)

\usepackage[ngerman]{babel}
\usepackage{geometry}
\geometry{a4paper, top=2.2cm, bottom=2.2cm, left=2.3cm, right=2.3cm}

\usepackage{fontspec}
% Fallback-Kette: Source Sans Pro -> Calibri -> Segoe UI -> Latin Modern Sans
% (Latin Modern ist immer mit MiKTeX vorhanden)
\IfFontExistsTF{Source Sans Pro}{%
  \setmainfont{Source Sans Pro}[Ligatures=TeX]%
}{\IfFontExistsTF{Calibri}{%
  \setmainfont{Calibri}[Ligatures=TeX]%
}{\IfFontExistsTF{Segoe UI}{%
  \setmainfont{Segoe UI}[Ligatures=TeX]%
}{%
  \setmainfont{Latin Modern Sans}[Ligatures=TeX]%
}}}
\IfFontExistsTF{Source Code Pro}{%
  \setmonofont{Source Code Pro}[Scale=0.92]%
}{\IfFontExistsTF{Consolas}{%
  \setmonofont{Consolas}[Scale=0.92]%
}{%
  \setmonofont{Latin Modern Mono}[Scale=0.92]%
}}

\usepackage{xcolor}
% Farbpalette: hell, freundlich, modern
\definecolor{bfwPrimary}{HTML}{0EA5A4}    % Petrol/Türkis - Primär
\definecolor{bfwPrimaryDark}{HTML}{0F766E} % dunkler Petrol für Headings
\definecolor{bfwAccent}{HTML}{F59E0B}     % Amber - Akzent
\definecolor{bfwSuccess}{HTML}{10B981}    % Grün - Erfolg / Lösung
\definecolor{bfwWarn}{HTML}{F97316}       % Orange - Achtung
\definecolor{bfwInk}{HTML}{0F172A}        % fast Schwarz
\definecolor{bfwInkSoft}{HTML}{475569}    % gedämpftes Grau
\definecolor{bfwBgSoft}{HTML}{F8FAFC}     % off-white
\definecolor{bfwBgTeal}{HTML}{ECFEFF}     % helles Türkis
\definecolor{bfwBgAmber}{HTML}{FEF3C7}    % helles Gelb
\definecolor{bfwBgRose}{HTML}{FFE4E6}     % helles Rosé für Eselsbrücken

\usepackage[colorlinks=true, linkcolor=bfwPrimaryDark, urlcolor=bfwPrimaryDark]{hyperref}
\usepackage{enumitem}
\setlist[itemize]{leftmargin=*, itemsep=2pt, topsep=2pt}
\setlist[enumerate]{leftmargin=*, itemsep=2pt, topsep=2pt}

\usepackage[most]{tcolorbox}
\usepackage{booktabs}
\usepackage{tikz}
\usetikzlibrary{decorations.pathmorphing, positioning, shapes.geometric, arrows.meta}

% Merkkästchen — wichtige Definition / Kernpunkt
\newtcolorbox{merksatz}[1][]{%
  enhanced, breakable,
  colback=bfwBgTeal,
  colframe=bfwPrimary,
  boxrule=0.6pt,
  arc=4pt,
  left=10pt, right=10pt, top=8pt, bottom=8pt,
  fonttitle=\bfseries\color{bfwPrimaryDark},
  title={\faLightbulb\ Merksatz},
  #1
}

% Eselsbrücke — Mnemoniks
\newtcolorbox{eselsbruecke}[1][]{%
  enhanced, breakable,
  colback=bfwBgRose,
  colframe=bfwAccent,
  boxrule=0.6pt,
  arc=4pt,
  left=10pt, right=10pt, top=8pt, bottom=8pt,
  fonttitle=\bfseries\color{bfwWarn},
  title={Eselsbrücke},
  #1
}

% Beispiel-Box
\newtcolorbox{beispiel}[1][]{%
  enhanced, breakable,
  colback=bfwBgSoft,
  colframe=bfwInkSoft,
  boxrule=0.4pt,
  arc=3pt,
  left=10pt, right=10pt, top=8pt, bottom=8pt,
  fonttitle=\bfseries\color{bfwInk},
  title={Beispiel},
  #1
}

% Lösungs-Box (für Aufgabenhefte)
\newtcolorbox{loesung}[1][]{%
  enhanced, breakable,
  colback=bfwBgAmber!40,
  colframe=bfwSuccess,
  boxrule=0.6pt,
  arc=3pt,
  left=10pt, right=10pt, top=8pt, bottom=8pt,
  fonttitle=\bfseries\color{bfwSuccess!70!black},
  title={Lösung},
  #1
}

% Glossar-Eintrag
\newcommand{\glossareintrag}[2]{%
  \par\noindent\textbf{\color{bfwPrimaryDark}#1} \enspace #2\par\smallskip
}

% Section-Stil — etwas farbiger als Standard
\usepackage{titlesec}
\titleformat{\section}
  {\Large\bfseries\color{bfwPrimaryDark}}
  {\thesection}{0.6em}{}
\titleformat{\subsection}
  {\large\bfseries\color{bfwPrimary}}
  {\thesubsection}{0.5em}{}
\titleformat{\subsubsection}
  {\normalsize\bfseries\color{bfwInk}}
  {\thesubsubsection}{0.4em}{}

% TikZ-Stil "skizzenhaft" — etwas weicher als Rough.js, aber stabil
\tikzset{
  roughbox/.style = {
    draw=bfwPrimaryDark, thick, fill=bfwBgTeal,
    rounded corners=4pt, inner sep=6pt
  },
  roughaccent/.style = {
    draw=bfwAccent, thick, fill=bfwBgAmber,
    rounded corners=4pt, inner sep=6pt
  },
  roughline/.style = {
    draw=bfwInkSoft, thick, -{Stealth[length=2.5mm]}
  }
}

% Bullet-Symbole (bewusst kein FontAwesome — vermeidet Font-Installations-Probleme)
\newcommand{\faLightbulb}{$\bullet$}
\newcommand{\faBookmark}{$\bullet$}

% Header/Footer
\usepackage{fancyhdr}
\pagestyle{fancy}
\fancyhf{}
\renewcommand{\headrulewidth}{0pt}
\fancyfoot[L]{\small\color{bfwInkSoft}\thepage}
\fancyfoot[R]{\small\color{bfwInkSoft} BFW M-064 Beschaffung im Büromanagement}


\title{\vspace*{-1.5cm}\Huge\bfseries\color{bfwPrimaryDark}Aufgabenheft Tag~4\\[0.4em]
  \large\color{bfwInkSoft}\textnormal{Ausführen von Bestellungen --- Übungen im IHK-Klausurformat}}
\author{\textsf{Beschaffung im Büromanagement}\\
\textsf{\small\copyright\ Can Siebert\,\textperiodcentered\,2026}}
\date{Selbstlernmaterial \textperiodcentered{} 15.05.2026}

\begin{document}
\maketitle
\thispagestyle{empty}

\begin{merksatz}[title={\bfseries So nutzen Sie dieses Aufgabenheft}]
Bearbeiten Sie alle Aufgaben zunächst \emph{ohne} in die Lösungen zu schauen.
Beachten Sie die \emph{Operatoren} und schreiben Sie Antworten in vollständigen Sätzen.
Lösungen ab Seite~\pageref{loesungen}.
\end{merksatz}

\tableofcontents
\vspace{1em}
\hrule

\section{Aufgaben}

\subsection*{Aufgabe~1 --- Bestellpunktverfahren rechnen (12 Punkte)}

\noindent\textbf{\color{bfwAccent}YouTube-Vorbereitung:}\enspace \href{https://www.youtube.com/watch?v=1AltCD8OMX8}{Bestellpunkt- und Bestellrhythmusverfahren mit Meldebestand} (\url{https://www.youtube.com/watch?v=1AltCD8OMX8})

\medskip
Eine Druckerei verbraucht pro Werktag durchschnittlich 50~Blatt Premium-Papier. Die
Lieferzeit eines neuen Auftrags beträgt 5~Werktage. Als Sicherheitsbestand werden
100~Blatt vorgehalten.

\begin{enumerate}[label=(\alph*)]
  \item \textbf{Erläutern} Sie den Unterschied zwischen \emph{Mindestbestand} und \emph{Meldebestand}. \hfill (3~P.)
  \item \textbf{Berechnen} Sie den Meldebestand für die Druckerei. \hfill (3~P.)
  \item \textbf{Berechnen} Sie den Höchstbestand, wenn die Bestellmenge 1\,000~Blatt beträgt. \hfill (3~P.)
  \item \textbf{Beurteilen} Sie: Welcher Vorteil bietet das Bestellpunktverfahren gegenüber dem Bestellrhythmus\-verfahren? \hfill (3~P.)
\end{enumerate}

\subsection*{Aufgabe~2 --- Bestellverfahren vergleichen (10 Punkte)}

\noindent\textbf{\color{bfwAccent}YouTube-Vorbereitung:}\enspace \href{https://www.youtube.com/watch?v=Hve3iP9C7o8}{Lernvideo zu Bestellverfahren} (\url{https://www.youtube.com/watch?v=Hve3iP9C7o8})

\medskip
\begin{enumerate}[label=(\alph*)]
  \item \textbf{Nennen} Sie die vier wichtigsten Bestellverfahren. \hfill (2~P.)
  \item \textbf{Vergleichen} Sie Bestellpunkt- und Bestellrhythmus\-verfahren in einer kleinen Tabelle (Trigger, Vor-/Nachteil). \hfill (4~P.)
  \item \textbf{Beurteilen} Sie: Welche Voraussetzungen muss ein Lieferant erfüllen, damit JIT-Beschaffung funktioniert? \hfill (4~P.)
\end{enumerate}

\subsection*{Aufgabe~3 --- JIT und Kanban (8 Punkte)}

\noindent\textbf{\color{bfwAccent}YouTube-Vorbereitung:}\enspace \href{https://www.youtube.com/watch?v=6FX0M1tQ6pw}{Just-in-Time Produktion einfach erklärt} (\url{https://www.youtube.com/watch?v=6FX0M1tQ6pw})

\medskip
\begin{enumerate}[label=(\alph*)]
  \item \textbf{Erläutern} Sie das Just-in-Time-Prinzip in eigenen Worten. \hfill (3~P.)
  \item \textbf{Beschreiben} Sie, wie das Kanban-Verfahren funktioniert. \hfill (3~P.)
  \item \textbf{Beurteilen} Sie: Welches Risiko hat die JIT-Beschaffung im Krisenfall (z.\,B.\ Halbleiterknappheit 2021/22)? \hfill (2~P.)
\end{enumerate}

\subsection*{Aufgabe~4 --- Optimale Bestellmenge (Andler-Formel) (15 Punkte)}

\noindent\textbf{\color{bfwAccent}YouTube-Vorbereitung:}\enspace \href{https://www.youtube.com/watch?v=W6L-IkVTlLY}{Andler-Formel: Optimale Bestellmenge berechnen} (\url{https://www.youtube.com/watch?v=W6L-IkVTlLY})

\medskip
Ein Bürodienstleister benötigt pro Jahr \textbf{12\,000~Aktenordner}. Die Bestellkosten
pro Bestellung betragen \textbf{40~€}, der Einkaufspreis pro Aktenordner \textbf{2{,}50~€}
und der Lagerkostensatz \textbf{20~\% p.\,a.}

\begin{enumerate}[label=(\alph*)]
  \item \textbf{Nennen} Sie die Andler-Formel. \hfill (2~P.)
  \item \textbf{Erläutern} Sie die zwei Kostenarten, deren Konflikt die Andler-Formel optimiert. \hfill (3~P.)
  \item \textbf{Berechnen} Sie die optimale Bestellmenge. \hfill (6~P.)
  \item \textbf{Berechnen} Sie, wie oft pro Jahr bestellt werden muss. \hfill (2~P.)
  \item \textbf{Beurteilen} Sie: Welche Annahmen der Andler-Formel sind in der Praxis oft unrealistisch? \hfill (2~P.)
\end{enumerate}

\subsection*{Aufgabe~5 --- Bestellüberwachung und Mahnung (10 Punkte)}

\noindent\textbf{\color{bfwAccent}YouTube-Vorbereitung:}\enspace \href{https://www.youtube.com/watch?v=3e8Q_Brd2_s}{Mahnung schreiben --- Inkasso --- gerichtliches Mahnverfahren} (\url{https://www.youtube.com/watch?v=3e8Q_Brd2_s})

\medskip
Sie haben am 5.~Mai 100~Aktenordner bei der \emph{Bürofix AG} bestellt. Liefertermin
war der 15.~Mai. Heute ist der 20.~Mai --- die Lieferung ist noch nicht angekommen.

\begin{enumerate}[label=(\alph*)]
  \item \textbf{Erläutern} Sie den Begriff \emph{Bestellüberwachung} und nennen Sie zwei Aspekte, die überwacht werden. \hfill (3~P.)
  \item \textbf{Nennen} Sie die drei kaufmännischen Mahnstufen. \hfill (3~P.)
  \item \textbf{Verfassen} Sie den Inhalt einer schriftlichen Erinnerung (1.~Mahnstufe) an die \emph{Bürofix AG}. \hfill (4~P.)
\end{enumerate}

\subsection*{Aufgabe~6 --- Nichtige und anfechtbare Rechtsgeschäfte (5 Punkte)}

\noindent\textbf{\color{bfwAccent}YouTube-Vorbereitung:}\enspace \href{https://www.youtube.com/watch?v=8huvlRCYCT0}{Rechtsgeschäfte einfach erklärt} (\url{https://www.youtube.com/watch?v=8huvlRCYCT0})

\medskip
\begin{enumerate}[label=(\alph*)]
  \item \textbf{Unterscheiden} Sie nichtige und anfechtbare Rechtsgeschäfte. \hfill (3~P.)
  \item \textbf{Nennen} Sie zwei Anfechtungsgründe nach BGB. \hfill (2~P.)
\end{enumerate}

\vspace{1em}
\noindent\textbf{Punkte gesamt: 60}

\newpage

\section{Lösungen}\label{loesungen}

\subsection*{Lösung Aufgabe~1}
\begin{loesung}
\textbf{(a)} \emph{Mindestbestand} ist der Sicherheitspuffer im Lager, der nie unterschritten
werden sollte. \emph{Meldebestand} ist der Bestellauslöse\-punkt: Sobald der Lagerbestand
den Meldebestand erreicht, wird neu bestellt. Der Meldebestand liegt höher als der
Mindestbestand --- er enthält den erwarteten Verbrauch während der Lieferzeit als Puffer.

\textbf{(b)} Meldebestand:
$$\text{Meldebestand} = \text{Mindestbestand} + (\text{Tagesverbrauch} \times \text{Lieferzeit}) = 100 + (50 \times 5) = \textbf{350 Blatt}$$

\textbf{(c)} Höchstbestand:
$$\text{Höchstbestand} = \text{Mindestbestand} + \text{Bestellmenge} = 100 + 1\,000 = \textbf{1\,100 Blatt}$$

\textbf{(d)} Vorteile des Bestellpunktverfahrens: Es reagiert auf den \emph{tatsächlichen}
Verbrauch (statt nur auf den Kalender). Dadurch kann der durchschnittliche Lagerbestand
niedriger gehalten werden, und auch bei schwankendem Verbrauch wird rechtzeitig nachbestellt.
Das Bestellrhythmus\-verfahren würde unabhängig vom realen Verbrauch zum festen Termin
bestellen --- bei plötzlich steigendem Bedarf wäre Engpass möglich.
\end{loesung}

\subsection*{Lösung Aufgabe~2}
\begin{loesung}
\textbf{(a)} Vier Bestellverfahren: \textbf{Bestellpunktverfahren, Bestellrhythmus\-verfahren,
Just-in-Time, Kanban}.

\textbf{(b)} Vergleich:

\begin{center}
\begin{tabular}{p{4cm} p{4.5cm} p{4.5cm}}
\toprule
\textbf{Kriterium} & \textbf{Bestellpunkt} & \textbf{Bestellrhythmus}\\
\midrule
Auslöser & Bestand erreicht Meldebestand & Festgelegter Zeitpunkt\\
Bestellmenge & Konstant & Variabel (auf Höchstbestand)\\
Vorteil & Reaktion auf reale Nachfrage & Planbarkeit, geringer Aufwand\\
Nachteil & Permanente Bestandsüberwachung & Bei Bedarfsspitzen Engpass\\
Geeignet für & A-Teile & B-/C-Teile\\
\bottomrule
\end{tabular}
\end{center}

\textbf{(c)} Voraussetzungen für JIT:
\begin{itemize}[itemsep=0pt]
  \item Sehr zuverlässige Lieferanten (hohe Termin- und Mengentreue)
  \item Kurze, planbare Lieferzeiten (oft geographische Nähe)
  \item Enge IT-Integration (EDI, gemeinsame Bestandsdaten)
  \item Disziplinierte Prozesse auf beiden Seiten
\end{itemize}
\end{loesung}

\subsection*{Lösung Aufgabe~3}
\begin{loesung}
\textbf{(a)} Just-in-Time bedeutet: Material wird genau zum Bedarfszeit\-punkt geliefert.
Lager praktisch null, Kapital­bindung minimal. Die Lieferung kommt fertigungs\-synchron,
oft mehrmals täglich, direkt an die Produktionslinie.

\textbf{(b)} Kanban (jap.\ Karte) ist ein Pull-Verfahren: Wenn der Verbraucher Material
benötigt, signalisiert er das mit einer Karte (oder elektronisch). Erst dann produziert/
liefert der Lieferant die signalisierte Menge. Volle Box kommt zurück --- Kreislauf
beginnt von vorne. Vorteil: Das System steuert sich selbst, ohne zentrale Planung.

\textbf{(c)} Größtes Risiko der JIT-Beschaffung: \emph{Lieferketten\-störungen} treffen
direkt die Produktion. Konkretes Beispiel: Halbleiter\-knappheit 2021/22 --- viele
Automobil\-werke mussten ihre Produktion tagelang oder wochenlang stoppen, weil
JIT-gelieferte Mikrochips ausblieben. Da kein Sicherheitsbestand vorhanden war,
führte der Lieferengpass sofort zum Produktions\-stopp mit Millionen­verlusten.
\end{loesung}

\subsection*{Lösung Aufgabe~4}
\begin{loesung}
\textbf{(a)} Andler-Formel:
\[
x_{\text{opt}} = \sqrt{\frac{2 \cdot \text{Jahresbedarf} \cdot \text{Bestellkosten}}{\text{Einkaufspreis} \cdot \text{Lagerkostensatz}}}
\]

\textbf{(b)} Die Andler-Formel optimiert den Konflikt zwischen \emph{Bestellkosten}
(Personal-, IT-, Frachtsockel-Kosten pro Bestellvorgang --- sinken pro Stück bei größerer
Bestellmenge) und \emph{Lagerkosten} (Lagermiete, Energie, Versicherung, Kapital­bindung
--- steigen mit dem Lagerbestand). Im Optimum ist die Summe beider Kostenarten am
niedrigsten.

\textbf{(c)} Berechnung:
\[
x_{\text{opt}} = \sqrt{\frac{2 \cdot 12\,000 \cdot 40}{2{,}50 \cdot 0{,}20}}
= \sqrt{\frac{960\,000}{0{,}50}}
= \sqrt{1\,920\,000}
\approx \textbf{1\,386 Stück}
\]

\textbf{(d)} Bestellhäufigkeit:
$$12\,000 \div 1\,386 \approx \textbf{8{,}7 Bestellungen pro Jahr} \approx \text{alle 6 Wochen.}$$

\textbf{(e)} Praxisferne Annahmen:
\begin{itemize}[itemsep=0pt]
  \item Konstanter Verbrauch über das Jahr (in Wirklichkeit: Saisonschwankungen)
  \item Konstanter Einkaufspreis --- ignoriert Mengenrabatte!
  \item Sofortige Lieferung ohne Mindestbestellmengen
\end{itemize}
\end{loesung}

\subsection*{Lösung Aufgabe~5}
\begin{loesung}
\textbf{(a)} \emph{Bestellüberwachung} = Kontrolle, ob die Lieferung termin- und
mengentreu eingeht. Sie startet mit dem Versenden der Bestellung und endet mit dem
ordnungs\-gemäßen Wareneingang. Zwei Aspekte: \textbf{Termintreue} (kommt die Lieferung
pünktlich?) und \textbf{Mengentreue} (stimmt die Menge?). Drittens auch: Qualität ---
das ist aber typisch beim Wareneingang (Tag~5).

\textbf{(b)} Drei Mahnstufen: 1.~\textbf{Erinnerung} (freundlicher Hinweis, neue Frist),
2.~\textbf{Mahnung} (klarer Verzugs\-hinweis, Schadensersatz angedroht), 3.~\textbf{letzte
Mahnung} (Androhung von Rücktritt, Schadensersatz und gerichtlichem Mahnverfahren).

\textbf{(c)} Musterhafte Erinnerung (1.~Mahnstufe):

\begin{quote}\small
Schreibwaren-Vertrieb GmbH, Berlin \\
20.~Mai 2026

\medskip
Bürofix AG, München

\medskip
\textbf{Erinnerung Bestellung Nr.\ 2026-0512-001 --- 100 Aktenordner}

\medskip
Sehr geehrte Damen und Herren,

\medskip
am 5.~Mai 2026 haben wir bei Ihnen 100~Aktenordner DIN~A4 (blau, breit) zum Liefertermin
\textbf{15.~Mai 2026} bestellt. Bis heute ist die Lieferung leider nicht eingetroffen.

\medskip
Wir bitten Sie höflich, die Lieferung bis spätestens \textbf{27.~Mai 2026} zu veranlassen.
Sollte uns die Ware bis dahin nicht erreicht haben, müssten wir formell mahnen
und uns weitere Schritte vorbehalten.

\medskip
Mit freundlichen Grüßen \\
Ihr Name --- Sachbearbeitung Einkauf
\end{quote}
\end{loesung}

\subsection*{Lösung Aufgabe~6}
\begin{loesung}
\textbf{(a)} Unterschied:
\begin{itemize}[itemsep=0pt]
  \item \emph{Nichtig}: von Anfang an unwirksam, ohne dass es einer Erklärung bedarf. Z.\,B.\ Geschäftsunfähigkeit (§\,105), Verstoß gegen Verbotsgesetz (§\,134), Sittenwidrigkeit (§\,138).
  \item \emph{Anfechtbar}: zunächst wirksam, wird aber durch eine Anfechtungs\-erklärung des Berechtigten rückwirkend nichtig (§\,142 BGB). Erfordert aktives Handeln.
\end{itemize}

\textbf{(b)} Zwei Anfechtungsgründe:
\begin{itemize}[itemsep=0pt]
  \item \emph{Inhalts- oder Erklärungs\-irrtum} (§\,119 BGB) --- Versprecher, falsches Verständnis
  \item \emph{Arglistige Täuschung oder Drohung} (§\,123 BGB) --- bewusste Irreführung oder Druck
\end{itemize}
\end{loesung}

\end{document}
